\section{Replicação Cruzada}
\label{sec:replicacao}

Em conformidade com as diretrizes do estudo, a replicação ocorreu nas duas direções sobre
uma subamostra determinística de 30 repositórios sorteada com semente fixa e documentada,
seguindo apenas o \texttt{README} publicado de cada lado.

\subsection{Como Replicadores}

Executamos o \emph{pipeline} do grupo sorteado no commit \texttt{73b739f}, em macOS 27.0.1
(arm64) com Python 3.14.0, em ambiente virtual novo e sem nenhum cache prévio. A subamostra
foi sorteada sobre o conjunto ordenado e deduplicado dos 300 repositórios do dataset
original deles, com o \emph{hash} SHA-256 do CSV de origem registrado junto da lista, de
modo que o sorteio seja refazível por qualquer pessoa --- inclusive pelo grupo autor ao
responder as Issues.

\textbf{O estudo foi executável}, com uma ressalva de escopo. Clone, instalação das nove
dependências e o modo de estimativa de consumo de API funcionaram sem atrito, somando menos
de dois minutos. A execução completa levou \textbf{42 min 20 s}, dos quais 38 min 38 s
(91\%) foram espera da API, consumindo cerca de 1.676 requisições para 30 repositórios ---
aproximadamente 2,5 vezes a projeção linear da própria estimativa do \emph{pipeline}. Todos
os artefatos prometidos foram produzidos e a suíte de testes do grupo passou integralmente
(159 testes, 77\% de cobertura).

A ressalva é relevante para a própria viabilidade da replicação: \textbf{não existe forma,
documentada ou não, de executar aquele \emph{pipeline} sobre um conjunto dado de
repositórios}. A etapa de seleção sempre refaz a busca por faixas de estrelas, que depende
da contagem de estrelas no instante da execução e, portanto, não é estável entre execuções.
Para rodar sobre a subamostra foi necessário ler o código-fonte, descobrir que as etapas
seguintes consomem um CSV intermediário e escrevê-lo manualmente --- isto é, sair do
\texttt{README}. Seguindo estritamente a documentação, a replicação cruzada como o estudo a
define não é executável naquele \emph{pipeline}.

Quanto aos valores, a reprodução foi praticamente exata: a diferença relativa mediana foi de
\textbf{0,0\% nas sete métricas} comparadas, com concordância de categoria de 100\% em
frequência, lead time e CFR, e de 91,3\% na recuperação ($\kappa$ ponderado de 0,907). As
divergências concentram-se em repositórios com poucos episódios de falha, em que um único
episódio a mais ou a menos desloca a mediana; o maior desvio relativo observado na
recuperação chega a 945\%, em um repositório com baixa contagem de episódios. Abrimos quatro
Issues no repositório do grupo autor: uma de execução e documentação, referente à ausência de
parâmetro para a subamostra, e três de documentação --- links apontando para o disco local de
um dos autores, instalação sem ambiente virtual e contagem de colunas desatualizada.

\subsection{Como Autores}

Recebemos três Issues do grupo replicador, todas procedentes e todas corrigidas antes do
prazo, e todas registradas no GitHub Projects do grupo.

A primeira relatava que o comando de análise falhava em um clone limpo, porque os CSVs do
estudo não estavam versionados na \emph{branch} principal --- existiam apenas em uma
\emph{branch} seis commits à frente, nunca integrada. Corrigimos publicando os três CSVs, de
modo que a análise passou a ser reexecutável a partir de um clone limpo, sem token e sem
refazer a coleta.

A segunda apontava a ausência de um parâmetro para executar sobre a subamostra --- exatamente
a mesma lacuna que havíamos encontrado no \emph{pipeline} deles. Acrescentamos uma opção que
recebe a lista de repositórios e substitui a busca por estrelas, preservando os critérios de
inclusão: um repositório da lista que não os atenda continua aparecendo no funil como
excluído, com o motivo registrado.

A terceira relatava que um repositório mudara de categoria geral entre o dataset publicado e
a reexecução. A investigação mostrou que \textbf{não havia erro de cálculo em nenhum dos dois
lados}: o repositório está 1,87 h acima do limite de 168 h entre \emph{Medium} e \emph{Low} no
tempo de recuperação --- 1,10\% do valor --- e, como a categoria geral é a mediana das quatro
notas arredondada para baixo, mover uma única métrica um degrau basta para alterá-la. Essa
Issue motivou a análise de fragilidade de fronteira relatada na Seção \ref{sec:resultados} e
a ameaça à validade de conclusão correspondente.

\subsection{O que a Replicação Mudou no Estudo}

Além das três correções, a tentativa de reproduzir a comparação revelou um defeito que
\emph{nenhuma} das Issues havia apontado: o módulo de análise esperava nomes de coluna que o
\emph{pipeline} de coleta não escreve. Como coluna ausente é tratada como valor ausente por
contrato, a análise lia o CSV real \textbf{sem erro algum e devolvia estatísticas vazias}. A
falha ruidosa relatada pelo grupo replicador encobria uma falha silenciosa; corrigida a
primeira, a segunda tornou-se visível. É o argumento mais concreto que este trabalho produziu
em favor da replicação cruzada como prática.
